\documentclass[10pt,a4paper]{amsart}
\usepackage{amsmath,amssymb,amsthm,graphicx}
% Retained minimal preamble for the Stacks Project chapter more-morphisms.tex (More on Morphisms).
% Verbatim excerpt of the chapter preamble: Xy-pic setup (native preamble lines 14-18) and the native
% theorem declarations with the mathematical macros (native preamble lines 154-217). Copied mechanically;
% nothing added, removed or redefined. The wrapper class and the amsmath/amssymb/amsthm/graphicx base
% packages are supplied by the assembler, so the class/externaldocument part of the native preamble is omitted.
% For commutative diagrams we use Xy-pic
\usepackage[all]{xy}

% We use 2cell for 2-commutative diagrams.
\xyoption{2cell}
% Theorem environments.
%
\theoremstyle{plain}
\newtheorem{theorem}[subsection]{Theorem}
\newtheorem{proposition}[subsection]{Proposition}
\newtheorem{lemma}[subsection]{Lemma}

\theoremstyle{definition}
\newtheorem{definition}[subsection]{Definition}
\newtheorem{example}[subsection]{Example}
\newtheorem{exercise}[subsection]{Exercise}
\newtheorem{situation}[subsection]{Situation}

\theoremstyle{remark}
\newtheorem{remark}[subsection]{Remark}
\newtheorem{remarks}[subsection]{Remarks}

\numberwithin{equation}{subsection}

% Macros
%
\def\lim{\mathop{\mathrm{lim}}\nolimits}
\def\colim{\mathop{\mathrm{colim}}\nolimits}
\def\Spec{\mathop{\mathrm{Spec}}}
\def\Hom{\mathop{\mathrm{Hom}}\nolimits}
\def\Ext{\mathop{\mathrm{Ext}}\nolimits}
\def\SheafHom{\mathop{\mathcal{H}\!\mathit{om}}\nolimits}
\def\SheafExt{\mathop{\mathcal{E}\!\mathit{xt}}\nolimits}
\def\Sch{\mathit{Sch}}
\def\Mor{\mathop{\mathrm{Mor}}\nolimits}
\def\Ob{\mathop{\mathrm{Ob}}\nolimits}
\def\Sh{\mathop{\mathit{Sh}}\nolimits}
\def\NL{\mathop{N\!L}\nolimits}
\def\CH{\mathop{\mathrm{CH}}\nolimits}
\def\proetale{{pro\text{-}\acute{e}tale}}
\def\etale{{\acute{e}tale}}
\def\QCoh{\mathit{QCoh}}
\def\Ker{\mathop{\mathrm{Ker}}}
\def\Im{\mathop{\mathrm{Im}}}
\def\Coker{\mathop{\mathrm{Coker}}}
\def\Coim{\mathop{\mathrm{Coim}}}

% Boxtimes
%
\DeclareMathSymbol{\boxtimes}{\mathbin}{AMSa}{"02}

%
% Macros for moduli stacks/spaces
%
\def\QCohstack{\mathcal{QC}\!\mathit{oh}}
\def\Cohstack{\mathcal{C}\!\mathit{oh}}
\def\Spacesstack{\mathcal{S}\!\mathit{paces}}
\def\Quotfunctor{\mathrm{Quot}}
\def\Hilbfunctor{\mathrm{Hilb}}
\def\Curvesstack{\mathcal{C}\!\mathit{urves}}
\def\Polarizedstack{\mathcal{P}\!\mathit{olarized}}
\def\Complexesstack{\mathcal{C}\!\mathit{omplexes}}
% \Pic is the operator that assigns to X its picard group, usage \Pic(X)
% \Picardstack_{X/B} denotes the Picard stack of X over B
% \Picardfunctor_{X/B} denotes the Picard functor of X over B
\def\Pic{\mathop{\mathrm{Pic}}\nolimits}
\def\Picardstack{\mathcal{P}\!\mathit{ic}}
\def\Picardfunctor{\mathrm{Pic}}
\def\Deformationcategory{\mathcal{D}\!\mathit{ef}}
\newenvironment{reference}{}{}
\newenvironment{slogan}{}{}
\newenvironment{history}{}{}
\usepackage{cleveref}
\crefname{originallemma}{original lemma}{original lemmas}
\Crefname{originallemma}{Original lemma}{Original lemmas}
\makeatletter
\newlabel{divisors-lemma-blow-up-pullback-effective-Cartier}{{0809}{source}{}{sourcepointer-1}{}}
\newlabel{divisors-lemma-blow-up-pullback-effective-Cartier@cref}{{[originallemma][0][]0809}{[1][0][]source}{}{}{}}
\makeatother
\crefname{originallemma}{original lemma}{original lemmas}
\Crefname{originallemma}{Original lemma}{Original lemmas}
\makeatletter
\newlabel{divisors-lemma-blowing-up-gives-effective-Cartier-divisor}{{02OS}{source}{}{sourcepointer-2}{}}
\newlabel{divisors-lemma-blowing-up-gives-effective-Cartier-divisor@cref}{{[originallemma][0][]02OS}{[1][0][]source}{}{}{}}
\makeatother
\crefname{originallemma}{original lemma}{original lemmas}
\Crefname{originallemma}{Original lemma}{Original lemmas}
\makeatletter
\newlabel{divisors-lemma-blowing-up-projective}{{02NS}{source}{}{sourcepointer-3}{}}
\newlabel{divisors-lemma-blowing-up-projective@cref}{{[originallemma][0][]02NS}{[1][0][]source}{}{}{}}
\makeatother
\crefname{originallemma}{original lemma}{original lemmas}
\Crefname{originallemma}{Original lemma}{Original lemmas}
\makeatletter
\newlabel{divisors-lemma-blowing-up-two-ideals}{{080A}{source}{}{sourcepointer-4}{}}
\newlabel{divisors-lemma-blowing-up-two-ideals@cref}{{[originallemma][0][]080A}{[1][0][]source}{}{}{}}
\makeatother
\crefname{originallemma}{original lemma}{original lemmas}
\Crefname{originallemma}{Original lemma}{Original lemmas}
\makeatletter
\newlabel{lemma-get-ample-family}{{0GTS}{source}{}{sourcepointer-5}{}}
\newlabel{lemma-get-ample-family@cref}{{[originallemma][0][]0GTS}{[1][0][]source}{}{}{}}
\makeatother
\crefname{originallemma}{original lemma}{original lemmas}
\Crefname{originallemma}{Original lemma}{Original lemmas}
\makeatletter
\newlabel{properties-lemma-point-and-maximal-points-affine}{{01ZX}{source}{}{sourcepointer-6}{}}
\newlabel{properties-lemma-point-and-maximal-points-affine@cref}{{[originallemma][0][]01ZX}{[1][0][]source}{}{}{}}
\makeatother
\crefname{originallemma}{original lemma}{original lemmas}
\Crefname{originallemma}{Original lemma}{Original lemmas}
\makeatletter
\newlabel{properties-lemma-quasi-coherent-finite-type-ideals}{{01PH}{source}{}{sourcepointer-7}{}}
\newlabel{properties-lemma-quasi-coherent-finite-type-ideals@cref}{{[originallemma][0][]01PH}{[1][0][]source}{}{}{}}
\makeatother
\title{Blowing up a finite type ideal near a curve point where the restriction is ample}
\author{The Stacks Project Authors}
\date{stacks-project:more-morphisms.tex:local-snapshot-20261011; GNU FDL 1.2 or later}
\begin{document}
\maketitle
\noindent\textit{Editorial note.} Statement span is lines 23391-23400 of sources/stacks-more-morphisms/pinned-original-source.tex; proof span is lines 23402-23438, the complete author proof (35 body lines).\par\medskip
\noindent\textit{Editorial note.} Navigation stubs: 7 cited-but-not-included labels are supplied as non-load-bearing pointers with their pinned tags where the snapshot tag index has them.\par\medskip
\noindent\textit{Editorial note.} extra\_packages is empty on purpose: the retained preamble already loads xy with options and a second declaration clashes in the assembler wrapper.\par\medskip
\noindent\textit{Editorial note.} Text is retained verbatim, including the author's own typos and placeholders.\par\medskip
\noindent\textit{Editorial note.} Terminal hub of the chapter's blow-up series; the only \textbackslash{}cite in the collected spans, satisfied by a card-local bibliography.tex (Gross-thesis entry from the snapshot my.bib).\par\medskip
\noindent\textit{Editorial note.} The cited proposition's number appears verbatim in the retained statement, so the reference list is load-bearing for a reader; entry transcribed exactly from my.bib lines 2654-2660.\par\medskip
\noindent\textit{Allowed standard prerequisites (editorial).}\begin{itemize}
\item Explicit prerequisite boundary. Gross, Vector bundles as generators on schemes and stacks (Proposition 1.3.1, an ampleness criterion on blow-ups checked formally along the centre), the chapter's preceding blow-up/slice lemmas used inside the proof, and standard ampleness machinery on the generic fibre.
\end{itemize}\medskip
\ifdefined\hypertarget\hypertarget{sourcepointer-1}{}\fi
\noindent\textit{Source reference.} Original lemma 0809: sources/stacks-more-morphisms/pinned-original-source.tex, source lines 23438--23438. Cited by the retained statement or proof but not part of the collected spans; supplied only as a navigation pointer. The argument does not depend on its content.\par\medskip
\ifdefined\hypertarget\hypertarget{sourcepointer-2}{}\fi
\noindent\textit{Source reference.} Original lemma 02OS: sources/stacks-more-morphisms/pinned-original-source.tex, source lines 23438--23438. Cited by the retained statement or proof but not part of the collected spans; supplied only as a navigation pointer. The argument does not depend on its content.\par\medskip
\ifdefined\hypertarget\hypertarget{sourcepointer-3}{}\fi
\noindent\textit{Source reference.} Original lemma 02NS: sources/stacks-more-morphisms/pinned-original-source.tex, source lines 23438--23438. Cited by the retained statement or proof but not part of the collected spans; supplied only as a navigation pointer. The argument does not depend on its content.\par\medskip
\ifdefined\hypertarget\hypertarget{sourcepointer-4}{}\fi
\noindent\textit{Source reference.} Original lemma 080A: sources/stacks-more-morphisms/pinned-original-source.tex, source lines 23438--23438. Cited by the retained statement or proof but not part of the collected spans; supplied only as a navigation pointer. The argument does not depend on its content.\par\medskip
\ifdefined\hypertarget\hypertarget{sourcepointer-5}{}\fi
\noindent\textit{Source reference.} Original lemma 0GTS: sources/stacks-more-morphisms/pinned-original-source.tex, source lines 23361--23367. Cited by the retained statement or proof but not part of the collected spans; supplied only as a navigation pointer. The argument does not depend on its content.\par\medskip
\ifdefined\hypertarget\hypertarget{sourcepointer-6}{}\fi
\noindent\textit{Source reference.} Original lemma 01ZX: sources/stacks-more-morphisms/pinned-original-source.tex, source lines 23438--23438. Cited by the retained statement or proof but not part of the collected spans; supplied only as a navigation pointer. The argument does not depend on its content.\par\medskip
\ifdefined\hypertarget\hypertarget{sourcepointer-7}{}\fi
\noindent\textit{Source reference.} Original lemma 01PH: sources/stacks-more-morphisms/pinned-original-source.tex, source lines 23438--23438. Cited by the retained statement or proof but not part of the collected spans; supplied only as a navigation pointer. The argument does not depend on its content.\par\medskip
\begin{lemma}
\label{lemma-blow-up-ample-family}
\begin{reference}
\cite[Proposition 1.3.1]{Gross-thesis}
\end{reference}
Let $X$ be a quasi-compact and quasi-separated scheme with finitely
many irreducible components. There exists a quasi-compact dense open
$U \subset X$ and a $U$-admissible blowing up $X' \to X$ such that the
scheme $X'$ has an ample family of invertible modules.
\end{lemma}

\begin{proof}
Let $\eta_1, \ldots, \eta_n \in X$ be the generic points of the
irreducible components of $X$. By
Properties, Lemma \ref{properties-lemma-point-and-maximal-points-affine}
and the fact that $X$ is quasi-compact we can find a finite affine
open covering $X = U_1 \cup \ldots \cup U_m$ such that
each $U_i$ contains $\eta_1, \ldots, \eta_n$. In particular
the quasi-compact open subset $U = U_1 \cap \ldots \cap U_m$ is dense in $X$.
Let $\mathcal{I}_i \subset \mathcal{O}_X$ be a finite type quasi-coherent
ideal sheaf such that $U_i = X \setminus Z_i$ where $Z_i = V(\mathcal{I}_i)$,
see Properties, Lemma \ref{properties-lemma-quasi-coherent-finite-type-ideals}.
Let
$$
f : X' \longrightarrow X
$$
be the blowing up of $X$ in the ideal sheaf
$\mathcal{I} = \mathcal{I}_1 \cdots \mathcal{I}_m$.
Note that $f$ is a $U$-admissible blowing up as $V(\mathcal{I})$
is (set theoretically) the union of the $Z_i$ which are disjoint from $U$.
Also, $f$ is a projective morphism and
$\mathcal{O}_{X'}(1)$ is $f$-relatively ample, see
Divisors, Lemma \ref{divisors-lemma-blowing-up-projective}.
By Divisors, Lemma \ref{divisors-lemma-blowing-up-two-ideals}
for each $i$ the morphism $f'$ factors as $X' \to X'_i \to X$
where $X'_i \to X$ is the blowing up in $\mathcal{I}_i$
and $X' \to X'_i$ is another blowing up (namely in the pullback
of the products of the ideals $\mathcal{I}_j$ omitting $\mathcal{I}_i$).
It follows from this that $D_i = f^{-1}(Z_i) \subset X'$ is an effective
Cartier divisor, see
Divisors, Lemmas \ref{divisors-lemma-blow-up-pullback-effective-Cartier} and
\ref{divisors-lemma-blowing-up-gives-effective-Cartier-divisor}.
We have $X' \setminus D_i = f^{-1}(U_i)$. As $\mathcal{O}_{X'}(1)$ is
$f$-ample, the restriction of
$\mathcal{O}_{X'}(1)$ to $X' \setminus D_i$ is ample. It follows from
Lemma \ref{lemma-get-ample-family}
that $X'$ has an ample family of invertible modules.
\end{proof}

\begin{thebibliography}{1}
\bibitem{Gross-thesis}
Gross, Philipp.
\newblock {\em Vector bundles as generators on schemes and stacks}.
\newblock D{\"u}sseldorf, Univ., Diss., 2010.
\end{thebibliography}
\end{document}
